\section{Requirements and Threat Model}
\label{sec:requirements}

\subsection{Functional requirements}

The requirements below were derived from the three problems in
Section~\ref{sec:introduction} and refined against the codebase. Each is labelled so that
the design (Section~\ref{sec:design}) and the evaluation (Section~\ref{sec:evaluation})
can refer back to it.

\begin{description}[style=nextline,leftmargin=1.2em]
  \item[FR1 --- Bounded context under flood.] The notification block injected into the
        receiver's prompt must grow sub-linearly in the number of inbound messages from a
        peer; repeated events of the same kind from the same peer must collapse into one
        counted line.
  \item[FR2 --- Escalations always visible.] A pending approval, a failed connection or a
        paid priority message must be rendered regardless of how much chatter arrived
        before it, and regardless of the rendering cap.
  \item[FR3 --- Attention is measured.] The receiver must record, per peer and per day,
        how many notification lines were rendered on its behalf, how many estimated tokens
        those lines injected, how many escalations were surfaced and how many events were
        suppressed past the cap. The record must be readable without starting the
        transport.
  \item[FR4 --- Prices are discoverable.] An agent must be able to publish attention prices
        per tier in the same registration file that already publishes its endpoint, so that
        a sender can learn the cost before sending, and preview it with a dry run.
  \item[FR5 --- Enforcement is opt-in and machine-readable.] With enforcement enabled, an
        inbound \code{message/send} from a sender that holds no active \code{message/send}
        grant and carries no valid stamp must be rejected \emph{before} it reaches the
        conversation log or the notification pipeline, with a JSON-RPC error whose
        \code{data} carries the price list and, when applicable, the reason the stamp failed.
        With enforcement disabled (the default) behaviour must be unchanged.
  \item[FR6 --- Grants are free stamps.] A peer that holds a \code{message/send} grant is
        never charged for standard delivery; the social trust encoded in a grant takes
        precedence over the economic mechanism.
  \item[FR7 --- Pay once, spend many.] A sender must be able to prepay a credit at a receiver
        with one on-chain payment and then send many messages, each carrying a stamp that is
        debited off-chain; no chain RPC may occur in the per-message path.
  \item[FR8 --- Portable proof of credit.] The receiver must return a signed certificate
        that binds the credit to the payment transaction, the two agent identities and the
        amount, verifiable by the holder against the receiver's known agent address.
  \item[FR9 --- Paid wake-ups.] A stamp whose value meets the receiver's published
        \code{priority} price must cause the message to render as an escalation, with a
        longer excerpt, and to wake the agent where the host supports it---including when
        the sender is a grant holder.
  \item[FR10 --- Metered free tier.] A \code{message/send} grant may carry a
        \code{notificationsPerWeek} constraint. When the holder exceeds it, further ordinary
        notifications in the window fold into one summary line; the messages themselves are
        still delivered and logged. Escalations are never folded.
  \item[FR11 --- Host parity.] Every mechanism must behave identically whether the host
        is the CLI, \code{tapd}, the OpenClaw plugin or the Hermes plugin, because all of
        them render the same drained batch.
\end{description}

\subsection{Non-functional requirements}

\begin{description}[style=nextline,leftmargin=1.2em]
  \item[NFR1 --- Backward compatibility.] A peer running a pre-project TAP version must be
        able to message and be messaged by a post-project peer. Stamps ride in message
        metadata that older peers ignore; enforcement defaults to off.
  \item[NFR2 --- Crash safety.] The ledgers are written by exactly one process (the
        transport owner) using the repository's existing \code{tmp + rename} pattern, and
        every wire interaction that moves money must be idempotent across a crash between
        payment, debit, journal write and event emission.
  \item[NFR3 --- Never lose a message.] Quota folding and coalescing may compress
        \emph{notifications}; they must never discard a \emph{message}. Any enforcement path
        that cannot answer (missing ledger, missing contact) must fail open.
  \item[NFR4 --- No new trusted third party.] The design must not introduce a service that
        both peers must trust in addition to the chain and XMTP.
  \item[NFR5 --- Repository conventions.] ESM, named exports, Biome lint/format, strict
        TypeScript with \code{noUnusedLocals}; every new CLI command documented in the
        canonical skill file; both E2E suites updated for behavioural changes.
\end{description}

\subsection{Threat model}
\label{sec:threat-model}

The adversary is a \emph{connected} peer: an agent that has completed the invite handshake
and is an active contact, but that may be buggy, over-eager or hostile. Unknown senders are
already rejected by the transport and are out of scope. The adversary can send arbitrary
\code{message/send} and \code{action/request} payloads at any rate, can replay its own
earlier messages, can present arbitrary stamps and can claim arbitrary transaction
hashes in a top-up. It cannot forge another agent's XMTP identity or OWS signatures. The
assets to protect are the receiver's LLM context, the receiver's wake budget, the
integrity of the postage ledger (no double spend, no minting) and the receiver's signing
key (no signing amplification). Section~\ref{sec:security} maps each threat to a
mitigation.
